\section{本章小结}
\label{sec:measure-summary}

本章从“集合有多大”出发。第~\ref{sec:measure-size-additivity}节用计数、长度和Jordan测量说明，有限个区间的逼近无法处理有理点集这样的集合，测量规则必须对可数运算封闭。第~\ref{sec:measure-spaces-def}节把问题分成两层：$\sigma$代数规定哪些集合可以测量，测度用可数可加性规定怎样赋值；单调性、次可加性、单调连续性和零测例外都由这两条规定推出。第~\ref{sec:measure-construction}节回答测度从哪里来：外测度把可数覆盖的代价变成集合函数，Carathéodory条件筛出可加的集合，由此得到Lebesgue测度、Lebesgue–Stieltjes测度及一般的延拓定理；$\pi$–$\lambda$定理保证在生成族上一致的测度处处一致。Cantor集说明零测与可数是两回事。第~\ref{sec:measure-measurable-maps}节与第~\ref{sec:measure-pushforward}节转向映射：可测映射让输入端能够回答输出端提出的集合问题，推前测度用原像把输入质量汇集到输出集合上。

同一个Borel结构上可以放置长度、点质量以及它们的混合，因此本章后三节比较同一空间上的不同测度。绝对连续比较零测集是否包含，相互奇异询问质量能否分别集中在互不相交的可测区域；两者都只读取集合质量，相互奇异也不要求距离上的分离。有限测度作差得到有符号测度，Hahn分解找出正负区域，Jordan分解给出唯一的、相互奇异的正负部分，总变差测度$|\eta|$累计不受抵消影响的变化总量。Lebesgue分解则相对于一个参考测度，把$\sigma$有限测度唯一地分成绝对连续部分与奇异部分；在实轴上以Lebesgue测度为参考，奇异部分还能再分成点质量与奇异连续质量，Cantor测度给出后者的具体例子。三种分解分别回答：哪些集合被一同忽略，正负赋值落在哪里，以及相对于参考测度哪些质量能被看见。

进入第~\ref{chap:integration}章时，我们已经有了可测集合、测度、可测函数及其极限运算，简单函数也已定义为可测集合示性函数的有限线性组合；有符号测度的总变差也已就绪。第~\ref{chap:integration}章从简单函数的积分出发，构造简单函数逼近，建立一般积分及其收敛定理。第~\ref{chap:densities-change-of-measure}章在积分可用之后，研究一个测度能否表示为参考测度的加权、怎样计算这种权重并用于换测度。需要强调的是，本章完成的测度分解只把绝对连续质量分了出来，并没有证明它一定能由密度表示；那是Radon–Nikodym定理的任务。
